% !TeX root = ../papers/nc-p3.tex
\section{七次元対合成分の順序付き有理商と一般安定群}
\label{sec:rational-ordered-quartet}
\begin{theorem}\label{thm:rational-ordered-quartet}
定理\ref{thm:involution-quartet-component}の成分で、対合により第一点と第二点、
第三点と第四点が交換される順序付き四点列を、その表示の通り指定する。
この印付き成分は既約で、指定始点の同型類の関数体は七変数の純超越拡大である。
一般の無印窓の射影安定群は指定対合の生成する位数2の群であり、最小周期は4である。
\end{theorem}

\begin{proof}[二次曲線と倍率の消去]
節\ref{sec:involution-quartet-component}の $x,X,u,U,B,F,k$ を用いる。
整合条件の第一式から $U$ を消した第二式を $k^2=N(X)/D(X)$ と書く。
分子と分母はともに $X$ について一次で、その係数行列式は非零である。
従って $x,u,B,F,k$ を五つの独立座標にでき、$X,U$ はこれらの有理式になる。
この置換の逆は標本 $t$ でも定義される。
各自由二座標 $v_i,w_i$ の行列を
\[
 \begin{pmatrix}2v_i+a_i&2w_i+b_i\\2w_i+c_i&2v_i+d_i\end{pmatrix}
\]
と書き、$V_i=4v_i+a_i+d_i$, $W_i=4w_i+b_i+c_i$ とする。
その行列式の消滅は
\[
 (V_i+W_i)(V_i-W_i)=(a_i-d_i+b_i-c_i)(a_i-d_i-b_i+c_i)
\]
である。右辺は標本 $t$ の素数2657で順に $2218,1246$ となりともに単元である。
従って各曲線が開トーラスになり、五つの基底変数を含むこの稠密表示は既約である。
前節の $t$ での局所一致と、そこで全逆変換が存在することから、これが所定の成分を与える。

順序を固定した四点枠を保つ基底変更は位置別の対角行列である。
射影対合の行列を再び $(01)(23)$ にするため、各位置で両対内の倍率比は共通の符号となる。
隣接二位置の符号の積が負なら正4、負2の関係固有空間が交換されるため、四つの符号は全て等しい。
各位置スカラーを除く連結な倍率と全体倍率は $P=Q=A=E=l=1$ で一意に除かれた。
残った同時負符号は、この規格化に戻したとき各語に
$(-1)^{\sum_j (i_j\bmod 2)+\sum_j\lfloor i_j/2\rfloor}$ を掛ける。
従って $x,u,B,F$ を固定し $k$ を反転する。各ブロックの $v_i,a_i,d_i$ は偶、
$w_i,b_i,c_i$ は奇である。
\[
 s_i=\frac{V_i+W_i}{a_i-d_i+b_i-c_i},\qquad
 y_i=\frac{s_i-1}{k(s_i+1)}
\]
と置く。符号は $s_i\mapsto 1/s_i$ と作用するので $y_i$ は不変である。
逆に $k$ を一つ取れば $s_i=(1+ky_i)/(1-ky_i)$ から二座標が戻る。
よって所求の順序付き射影商の関数体は
$k_0(x,u,B,F,k^2,y_1,y_2)$（$k_0$ は基礎体）である。
\path{computations/verify_quartet_rational_chart.py} が一次分数式の次数・行列式、
二つの右辺と全係数の符号を検算する。
\end{proof}

\begin{proof}[印を外した自己同型]
$t$ で二十点が被約で、一般位置の五点から各位置の射影写像を一意に決められる。
ここで一般位置の五点とは、そのどの四点も独立であることをいう。
有限エタールに番号を付ければ、一般側の点配置間の射影写像は $t$ まで特殊化する。
実際、行列を射影空間に延ばし、階数が落ちた場合の核と像に入る点数を比べる。
$t$ では三点直線も五点平面もないので、階数 $r<4$ の行列では、
核の外の全点を相異なる像点に送れない。従って特殊化行列も可逆である。
隣接辺のグラフから一位置での番号の置換は全位置の置換を決定する。
系\ref{cor:klein-four-configurations}により、安定群は一般側で $t$ の四元群の部分群であり、
異なる位置の配置間には同型がない。従って一般周期は4である。

一般側で四元群全体が残ると仮定する。付値環上で二つの射影対合を線型に持ち上げ、
単元の平方根を加えて二乗を1とできる。交換子の符号は特殊側で正なので一般側でも正である。
四文字の多重度は各位置で独立な四つの文字のままであり、固有射影で局所的に文字枠を取れる。
隣接関係の色階数も固有空間の束の階数で一定、薄い全成分と可逆な交差ブロックは
$t$ で非零なので一般側でも非零である。
定理\ref{thm:klein-all-pairings}によればこの四元作用を持つ窓の無印像は次元5以下である。
これは前節の七次元性に反する。指定対合が残るので一般安定群は位数2となる。
\end{proof}
